% !TeX spellcheck = ru_RU
% !TEX root = vkr.tex

\section*{Введение}
\thispagestyle{withCompileDate}

Политика импортозамещения, подкрепленная в 2025–2026 годах субсидированием разработки программного обеспечения \cite{itGrant}, увеличивает потребность российской экономики в отечественных аппаратно-программных решениях. Санкционные ограничения и прекращение официальной поддержки предыдущих версий Windows, стимулирующее переход на Windows 11, способствуют росту доли операционной системы Linux на рынке \cite{osIncrese}, тогда как доля программного обеспечения в отрасли информационных технологий достигла 40\% в 2024 году при одновременном сокращении аппаратной составляющей. Эти изменения повышают спрос на разработчиков, владеющих фундаментальными принципами построения компьютерных систем — процессов, памяти, операционных систем — и способных реализовывать их на практике.

Одновременно активно развиваются технологические направления, требующие глубоких системных знаний: тяжёлые мультироторные дроны, робототехника, интернет вещей, процессоры для вычислительной техники, где критически важны понимание архитектуры набора инструкций (ISA), устройства центрального процессора (арифметико-логического устройства, регистров), взаимодействия с периферией и основ функционирования операционных систем. Рынок труда фиксирует значительный дефицит таких специалистов: в сфере информационных технологий открыто до миллиона вакансий, в микроэлектронике не хватает 5,3 тыс. инженеров (по данным 2024 г.), причем более 60\% позиций остаются незакрытыми из‑за пробелов в системном мышлении соискателей \cite{elStatistic}.

В российских университетах (БГТУ, СПбГУ и других) теория компьютерных систем в рамках дисциплин «Операционные системы» и «Архитектура ЭВМ» излагается подробно, однако материал подаётся фрагментарно. Курсы по архитектуре ЭВМ рассматривают модель фон Неймана, логические элементы, память и базовые компоненты процессоров; дисциплины по операционным системам — планировщики, прерывания и системные вызовы. В учебном процессе отсутствует наглядная демонстрация взаимосвязей между уровнями организации вычислительной системы: каким образом логические вентили образуют центральный процессор, как цикл операционной системы реализуется на уровне машинного кода, каким образом системные абстракции проявляются в ядре. Студенты приобретают практический опыт в симуляторах логических схем (Logisim, Digital), однако не переходят к более сложным и содержательным примерам — отладке программ, изучению минимальных операционных систем или виртуальных устройств. В результате выпускники владеют отдельными абстрактными концепциями, но затрудняются применять их для решения прикладных задач, что ограничивает их участие в проектной деятельности.

Открытые образовательные проекты — nandgame, nand2tetris, учебная ОС xv6, ранние версии Linux и GNU Mes — демонстрируют эффективный системный подход и представляют качественные практические примеры. Вместе с тем эти проекты недостаточно встроены в учебные программы, а начинающим разработчикам трудно самостоятельно выстроить на их основе целостный курс изучения материала.

Настоящая магистерская работа направлена на устранение указанного разрыва. Предлагается разработать учебный программный комплекс, объединяющий три ключевых компонента в единой воспроизводимой среде:

\begin{itemize}
  \item Виртуальный RISC‑V процессор на SystemVerilog с тактовой трассировкой и выводом временных диаграмм;
  \item Минимальную UNIX‑подобную операционную систему xv6, выполняющуюся непосредственно на виртуальном процессоре;
  \item Легковесный компилятор TinyCC, портированный для работы внутри xv6, позволяющий компилировать и выполнять пользовательские C-программы.
\end{itemize}

Комплекс реализует учебный процесс: обучающийся собирает операционную систему xv6 с помощью RISC‑V инструментария, загружает её в виртуальный процессор, затем внутри гостевой ОС запускает встроенный компилятор TinyCC, пишет программу на языке Си, компилирует её прямо в среде xv6 и наблюдает потактовое выполнение — стадия выполнения инструкции, обращения к памяти, обработку прерываний и системных вызовов — вплоть до вывода результата на виртуальный терминал. Тем самым в рамках одного лабораторного стенда связываются курсы «Архитектура ЭВМ», «Системное программное обеспечение» и «Операционные системы». Для подтверждения работоспособности комплекса планируется провести его тестирование на студенческих группах, а также подготовить комплект учебно-методических материалов и рекомендации по использованию смежных проектов. Программные компоненты комплекса реализуются на языках Verilog и Си и документируются с целью подтверждения возможности создания предложенных виртуальных устройств и их изучения в образовательном процессе.

